\section{从 Kernel 到系统：纯函数构件与良构边界}

上一章已经把 \texttt{kernels.cu} 的数值核心展开到 CUDA 执行结构上：
单边雅可比（one-sided Jacobi）把列对组织成无冲突 round，block 内规约（reduction）
计算局部 Gram 信息，kernel launch 边界承担 round、统计、旋转与后处理之间的全局同步。
这些内容回答的是一个局部而硬核的问题：\emph{如何让 GPU 正确执行 SVD 的核心变换}。

本章换一个视角。一个可运行、可维护、可复现实验的数值系统，还必须回答另一个问题：
\emph{这个 kernel 应该以什么形态嵌入系统}。如果 GPU kernel 直接认识文件格式、路径、
命令行参数、输出队列、线程池和构建目录，那么最昂贵、最难调试的一层就会被迫承担最多变化。
这不是高性能，这是把复杂度塞进最脆弱的位置。

本章的核心论点是：CUDA kernel 应当被设计成系统中的纯函数式构件
（pure-function-like component），而不是应用程序本身。严格地说，CUDA kernel 会写设备内存、
依赖 runtime 调度、消耗 GPU 资源，并不满足数学意义上的纯函数（pure function）。但是在系统边界上，
它应当表现得像一个确定的映射：
\[
    \operatorname{svd}_{\mathrm{cuda}}:
    (A,C)
    \longmapsto
    (U,\Sigma,V,r).
\]
其中 \(A\) 是主机侧行主序输入矩阵，\(C\) 是类型化计算配置，
\(U,\Sigma,V\) 是 SVD 输出，\(r\) 是当前 kernel 层能够报告的计算摘要。
在当前实现中，这个摘要具体表现为 \texttt{JacobiSvdResult::sweeps}；
更完整的端到端执行报告，例如 testcase 数、输出矩阵数、布局转置自动调优结果，
由应用层 \texttt{PipelineReport} 汇总，而不是塞进 kernel 的语义内部。

这个区分很重要。kernel 可以返回“我对这张矩阵迭代了多少 sweep”；pipeline 才能回答
“本次输入流一共处理了多少 testcase、写出了多少矩阵、运行时阈值是多少”。
前者是单矩阵数值事实，后者是端到端系统事实。把这两类事实混在一起，会让接口表面看起来更方便，
实际却更难测试、更难替换、更难解释。

\subsection{边界的必要性：高性能代码最怕职责膨胀}

GPU kernel 同时承受三类复杂性。

\begin{itemize}
    \item \textbf{数学复杂性}：单边雅可比需要维护列正交性、收敛阈值、奇异值提取、
          左右奇异向量关系以及 \(m\ge n\) 的 thin SVD 前置条件；
    \item \textbf{体系结构复杂性}：CUDA 的 thread、block、warp、shared memory、
          global memory、kernel launch 与 host-device 同步都会影响正确性和性能；
    \item \textbf{资源复杂性}：设备内存分配、主机到设备拷贝、设备到主机拷贝、
          错误码翻译、同步时机以及临时缓冲区生命周期必须被严格管理。
\end{itemize}

如果再让 kernel 层理解 \texttt{*.mat} 头部、文本矩阵空行、输出路径、CLI 默认值、
线程池窗口和 JSON 风格报告，它就会变成所有变化的交汇点。这样的设计在小 demo 里看似直接，
在真实演化中会迅速失控：文件格式变化影响 kernel，命令行变化影响 kernel，输出顺序变化影响 kernel，
最终最难验证的一层承担了最多责任。

良构边界（well-formed boundary）的目标不是把代码切成漂亮图层，而是降低变化传播
（change propagation）。理想情况下：
\[
    \Delta\text{CLI}
    \nrightarrow
    \Delta\text{Kernel},
    \qquad
    \Delta\text{File Format}
    \nrightarrow
    \Delta\text{Jacobi Rotation},
\]
\[
    \Delta\text{Output Ordering}
    \nrightarrow
    \Delta\text{CUDA Memory Layout},
    \qquad
    \Delta\text{Build Artifact}
    \nrightarrow
    \Delta\text{SVD Semantics}.
\]
这意味着用户接口、持久化格式、线程调度与构建方式都可以独立演化，
而不会倒灌进数值核心。对 CUDA 数值程序来说，这不是学院派洁癖，而是长期调试成本的生命线。

\subsection{轻量 DDD：领域不是业务，而是数值不变量}

领域驱动设计（Domain-Driven Design, DDD）常见于复杂业务系统，例如订单、账户与库存。
若机械套用，把每个类都命名成实体（Entity）、值对象（Value Object）或仓储（Repository），
这个项目会立刻沾上不必要的企业应用气味。这里真正有用的不是术语，而是 DDD 的朴素核心：
\emph{系统应围绕领域不变量组织，而不是围绕外部技术细节组织}。

在本项目中，“领域”（domain）不是业务规则，而是 SVD 的数值语义：
\[
    A \approx U\Sigma V^T,
    \qquad
    U^TU \approx I,
    \qquad
    V^TV \approx I,
\]
以及单边雅可比迭代中的局部不变量：
\[
    a_p^Ta_q \rightarrow 0,
    \qquad
    \lVert a_p\rVert_2,\lVert a_q\rVert_2
    \ \text{由正交旋转稳定更新}.
\]
这些关系才是系统真正要保护的核心。文件扩展名、命令行拼写、队列容量、CMake 生成器
只是让这个核心进入现实世界的外层机制。

因此，本项目采用的是轻量 DDD：只保留“领域核心、应用编排、基础设施、表示层”
这些边界思想，不强行引入完整战术模式。当前结构可以理解为：
\[
    \begin{array}{ccl}
        \text{Presentation}   & : & \texttt{main.cpp}\quad \text{把用户字符串编译成配置};       \\
        \text{Application}    & : & \texttt{pipeline.cpp}\quad \text{编排输入、计算、输出与并发}; \\
        \text{Domain Core}    & : & \texttt{kernels.cu}\quad \text{执行单矩阵 SVD 数值变换};  \\
        \text{Infrastructure} & : & \texttt{io.cu}\quad \text{处理矩阵流、字节、平台资源}.
    \end{array}
\]
注意这里使用 \texttt{io.cu}，因为当前 I/O 层已经包含 \texttt{cudaMallocHost} 管理的
页锁定主机缓冲区（pinned host buffer）。它仍然是基础设施层，而不是算法层；
文件、字节序、pinned staging 和文本兼容格式都属于数据进入系统的边界工作。

\subsection{Kernel 的公开契约：窄入口，类型化配置，主机侧结果}

当前 kernel 层向外暴露的核心入口是：
\[
    \begin{aligned}
        \texttt{one\_sided\_jacobi\_svd}:
         & \quad
        (\texttt{host\_input}, rows, columns, C) \\
         & \quad
        \longrightarrow \texttt{JacobiSvdResult}.
    \end{aligned}
\]
这个接口有三个值得保留的性质。

\paragraph{第一，输入是内存中的类型化矩阵，而不是外部世界。}
kernel 入口只接收主机侧行主序（row-major）数值数组、行列数和配置。
它不接收路径、不接收文件句柄、不接收网络序字节流，也不接收 \texttt{argv}。
坏文件必须在 I/O 层被拒绝，坏命令行必须在 CLI 层被拒绝；kernel 只处理已经成为矩阵的数据。

\paragraph{第二，配置显式表达数值与执行策略。}
\texttt{JacobiSvdConfig} 目前包含：
\[
    \epsilon,\quad
    \texttt{max\_sweeps},\quad
    \texttt{threads\_per\_block},
\]
以及布局转置相关策略：
\[
    \begin{gathered}
        \texttt{layout\_transpose\_mode},\quad
        \texttt{layout\_transpose\_min\_columns},\\
        \texttt{layout\_transpose\_min\_elements}
    \end{gathered}
\]
\[
    \begin{gathered}
        \texttt{layout\_transpose\_auto\_tune},\\
        \texttt{layout\_transpose\_benchmark\_repetitions},\\
        \texttt{layout\_transpose\_benchmark\_sweeps}.
    \end{gathered}
\]
这些字段并不全是同一种语义。\(\epsilon\) 与 \texttt{max\_sweeps} 更接近数值算法参数；
\texttt{threads\_per\_block} 与布局转置阈值更接近执行策略（execution policy）；
自动调优字段则描述“如何在本次运行前测量阈值”。当前把它们放在同一个结构体中是可接受的，
因为项目规模还小；但文档必须明确这种混合，以免未来把所有性能遥测和设备状态都继续塞进去。

\paragraph{第三，输出是主机侧值对象，而不是设备指针。}
\texttt{JacobiSvdResult} 包含：
\[
    rows,\quad columns,\quad u,\quad sigma,\quad v,\quad sweeps.
\]
其中 \(U\) 为 \(m\times n\) 行主序矩阵，\(\Sigma\) 为长度 \(n\) 的数组，
\(V\) 为 \(n\times n\) 行主序矩阵。设备指针、临时统计数组、旋转参数数组和布局工作区都不逃出
\texttt{kernels.cu} 的边界。外部世界拿到的是可移动、可测试、可序列化的 C++ 值。

这就是“纯函数式外观”的核心：内部可以有 \texttt{cudaMalloc}、\texttt{cudaMemcpy}、
kernel launch 和 \texttt{cudaDeviceSynchronize}；外部看到的是
\[
    (A,C)\mapsto(U,\Sigma,V,s).
\]
副作用被收束在窄边界内，接口语义因此接近数学函数。

\subsection{布局转置与自动调优：策略可以变化，契约不能变化}

上一章已经说明：外部契约保持行主序，但单边雅可比天然按列访问。
当前 kernel 层因此提供两条内部路径：

\begin{itemize}
    \item 直接路径：工作矩阵 \(A\) 保持行主序，列访问带有跨步；
    \item 布局转置路径：先在 GPU 上把 \(A\) 转为列主序工作区，再执行 sweep，收敛后转回行主序。
\end{itemize}

系统边界上最重要的不是“某条路径一定更快”，而是两条路径必须拥有同一公开契约。
无论 \texttt{LayoutTransposeMode} 是 \texttt{force\_enable}、\texttt{force\_disable}
还是 \texttt{auto\_select}，调用者最终都得到同形状、同行主序、同语义的 \(U,\Sigma,V\)。
布局是实现策略，不是外部 ABI。

自动调优（auto-tuning）进一步加强了这个边界意识。函数
\[
    \begin{aligned}
        \texttt{auto\_tune\_layout\_transpose\_threshold}:
         & \quad \texttt{JacobiSvdConfig}                               \\
         & \quad \longrightarrow \texttt{LayoutTransposeAutoTuneReport}
    \end{aligned}
\]
执行一组确定性微基准，返回推荐阈值、样本数与估计加速比。它不修改全局状态，也不偷偷改变用户配置。
当前 pipeline 的做法是：在正式读取输入流和提交线程池任务之前，若用户启用自动调优且模式为
\texttt{auto\_select}，就先运行一次调优，把推荐阈值写入本次运行的
\[
    C_{\mathrm{runtime}}.
\]
随后关闭该配置中的自动调优标志，避免 worker 线程重复执行微基准。

这个顺序很有品味。自动调优是运行前配置归一化，不是每个 testcase 自己执行的隐式副作用。
如果把调优放进工作任务内部，就会出现重复 benchmark、设备争用、报告不可解释等问题。
现在的结构则保持了清晰关系：
\[
    \text{user config}
    \xrightarrow{\text{optional pre-run tuning}}
    \text{runtime config}
    \xrightarrow{\text{kernel calls}}
    \text{results}.
\]
调优改变执行策略，但不改变 SVD 的外部语义。

\subsection{反向依赖是坏味道：kernel 不应该认识世界}

为了让边界足够锋利，需要明确列出 kernel 不应该知道的东西。

\begin{itemize}
    \item \textbf{不认识文件格式。} \texttt{kernels.cu} 不解析 \texttt{*.mat} 的网络字节序头部，
          也不解析 \texttt{*.txt} 中的空格、换行和空行；
    \item \textbf{不认识路径。} 输入路径、输出路径、父目录创建和覆盖策略都属于 CLI、pipeline 与 I/O；
    \item \textbf{不认识输出顺序。} \(U,\Sigma,V\) 如何被写为三张矩阵、乱序完成如何恢复顺序，
          是 pipeline 的职责；
    \item \textbf{不认识线程池。} 全局线程池（global thread pool）、future 窗口和重排序缓冲区
          （reorder buffer）服务于跨 testcase 调度，不属于单矩阵 SVD；
    \item \textbf{不认识用户意图。} \texttt{--epsilon}、\texttt{--max-sweeps}、
          \texttt{--layout-transpose} 在进入 kernel 前应已转化为类型化配置；
    \item \textbf{不认识构建系统。} \texttt{nvcc} driver mode、CMake CUDA language、
          输出 artifact 路径都不应影响运行时 SVD 语义；
    \item \textbf{不拥有长期系统资源。} 文件句柄、映射视图、输出线程、完成队列和 CLI 报告格式
          都不属于 kernel。
\end{itemize}

这组否定边界看似啰嗦，但它们能阻止架构坍缩。系统设计经常不是靠“应该做什么”保持清楚，
而是靠“绝不做什么”防止职责膨胀。对本项目而言，kernel 的尊严来自克制：
它不处理世界，它只处理矩阵。

\subsection{应用层：Pipeline 是用例协议，不是更大的 kernel}

\texttt{pipeline.cpp} 位于应用层（application layer）。它的职责不是实现 SVD，
而是组织一次完整用例（use case）：
\[
    \text{read testcase}
    \longrightarrow
    \text{compute SVD}
    \longrightarrow
    \text{write result}.
\]
但当前实现已经不再是简单的“主线程读、主线程算、写线程写”。新的 pipeline 包含：
\[
    \begin{gathered}
        \text{single-cursor dispatch}
        \longrightarrow
        \text{global thread pool}
        \longrightarrow
        \text{kernel stage} \\
        \longrightarrow
        \text{completion queue}
        \longrightarrow
        \text{reorder buffer}
        \longrightarrow
        \text{single writer}.
    \end{gathered}
\]

这套结构回答的是系统问题，而不是数学问题。不同 testcase 可以并发解码和计算；
完成顺序可以乱序；输出文件必须仍然按输入序写出
\[
    (U_0,\Sigma_0,V_0),
    (U_1,\Sigma_1,V_1),
    \ldots
\]
于是 pipeline 需要序号、future、完成队列、重排序缓冲区、背压（backpressure）和跨线程异常传播。
这些机制都不改变 Jacobi 旋转公式，但它们决定端到端程序是否可靠。

从边界角度看，pipeline 拥有的是执行顺序，而不是内部语义：

\begin{itemize}
    \item 它可以决定何时执行布局转置自动调优，但不实现转置 kernel；
    \item 它可以验证 testcase 的基本形状，但不实现 Jacobi 旋转；
    \item 它可以把 \texttt{JacobiSvdResult} 重组为三张 \texttt{io::Matrix}，
          但不改变 \(U,\Sigma,V\) 的数学含义；
    \item 它可以用重排序缓冲区恢复输出顺序，但不让 kernel 认识 testcase index。
\end{itemize}

这就是应用层聚合根（aggregate root）的意义：它把多个对象组合成一个端到端执行边界，
同时不偷走每一层自己的职责。pipeline 是用例协议，不是更大的 kernel。

\subsection{基础设施层：把脏世界封装在 I/O 边界}

\texttt{io.cu} 处理的是基础设施层（infrastructure layer）问题：文件、字节序、文本解析、
内存映射、平台 API、页锁定主机缓冲区和矩阵流 policy。这些问题非常重要，但它们不是 SVD 的领域语义。

当前 I/O 层提供两类入口。

\paragraph{第一，传统矩阵流。}
\[
    \texttt{MatInputStream},\quad
    \texttt{TxtInputStream},\quad
    \texttt{MatOutputStream},\quad
    \texttt{TxtOutputStream}.
\]
它们围绕 \texttt{io::Matrix} 工作：
\[
    \texttt{Matrix}=(rows,columns,values).
\]
这一接口适合简单路径、文本调试和兼容 API。policy-based design 让二进制格式与文本格式共享流控制，
而把具体编码策略隔离在 policy 中。

\paragraph{第二，高吞吐派发路径。}
\texttt{MatDispatchReader} 读取 \texttt{*.mat} 记录，并产生：
\[
    \texttt{MatDispatchTask}
    =
    (\texttt{sequence\_index},\texttt{rows},\texttt{columns},\texttt{buffer}).
\]
其中 \texttt{buffer} 是 \texttt{PinnedHostTaskBuffer}，由 \texttt{cudaMallocHost}/
\texttt{cudaFreeHost} 管理一块连续页锁定主机内存。输入区保存原始网络序 payload，
工作区为后续解析或 staging 预留空间。该路径让主线程保留单游标读取权，
同时把解码和计算移交给线程池任务。

这里需要保持清醒：pinned memory 仍然不是 device memory。I/O 层可以管理适合高速派发的主机缓冲，
但它不启动 kernel，不分配 \texttt{DeviceMatrix}，不管理 CUDA stream，也不决定布局转置策略。
它只是把外部字节转化为系统内部稳定任务或矩阵：
\[
    \text{external bytes}
    \longrightarrow
    \texttt{MatDispatchTask}
    \longrightarrow
    \texttt{io::Matrix}
    \longrightarrow
    \text{kernel input}.
\]
每个箭头都属于不同边界。把它们揉在一起，短期可能少几行代码，长期一定增加生命周期耦合。

\subsection{表示层：CLI 是意图编译器}

\texttt{main.cpp} 位于表示层（presentation layer）。它面对的是字符串：
\[
    argv=(a_0,a_1,\ldots,a_{n-1}).
\]
这些字符串本身没有领域意义。只有经过解析、类型转换、默认值归一化和跨字段验证后，
它们才变成应用层配置：
\[
    argv
    \longrightarrow
    \texttt{CliOptions}
    \longrightarrow
    \texttt{PipelineConfig}
    \longrightarrow
    \texttt{JacobiSvdConfig}.
\]

因此，CLI 可以被看成意图编译器（intent compiler）。用户表达“我要对这个文件做 SVD，
用这个收敛阈值，按这个格式写出结果”，CLI 把这句话降阶为 pipeline 能执行的配置。
它不应该直接读取矩阵文件，不应该启动 CUDA kernel，也不应该自己写 \(U,\Sigma,V\)。

这种边界带来一个朴素但强大的性质：用户接口可以变化，应用核心不必变化。
未来加入 \texttt{--dry-run}、\texttt{--device}、\texttt{--profile}
或 shell completion，
只要最终仍然生成稳定的 \texttt{PipelineConfig}，pipeline 与 kernel 就不需要知道用户使用了哪种表面语法。

\subsection{资源所有权：RAII 是边界纪律，不只是 C++ 技巧}

资源获取即初始化（Resource Acquisition Is Initialization, RAII）在本项目中不只是 C++ 写法偏好，
而是边界纪律。它回答一个非常实际的问题：谁负责释放什么，什么时候释放，失败路径是否泄漏。

当前系统中至少有这些资源所有者：

\begin{itemize}
    \item \texttt{DeviceMatrix} 拥有设备矩阵内存；
    \item 匿名命名空间中的 \texttt{DeviceBuffer<T>} 拥有临时设备数组；
    \item \texttt{MatReader} 与 \texttt{MatWriter} 拥有文件映射、句柄与映射视图；
    \item \texttt{PinnedHostTaskBuffer} 拥有页锁定主机内存；
    \item \texttt{ReorderBufferOutputStage} 拥有写线程、完成队列和关闭协议；
    \item 全局线程池拥有 worker 线程与任务队列。
\end{itemize}

每个资源都应有唯一所有者（unique owner）：
\[
    \operatorname{owner}(r)=o,
    \qquad
    \operatorname{lifetime}(r)\subseteq\operatorname{lifetime}(o).
\]
这条关系能消除许多 CUDA/C++ 混合程序中的经典错误：重复释放、早释放、异常路径泄漏、
线程未 join、映射视图悬垂、device pointer 漂移到上层。

RAII 与纯函数式外观并不矛盾。恰恰相反，正因为内部副作用由 RAII 对象管理，
外部接口才可以表现得更接近函数。资源纪律越清楚，函数式边界越可信。

\subsection{错误边界：把失败翻译成上层能理解的语言}

系统中的错误来源可以分成几类：
\[
    E =
    E_{\mathrm{cli}}
    \cup
    E_{\mathrm{io}}
    \cup
    E_{\mathrm{kernel}}
    \cup
    E_{\mathrm{pipeline}}
    \cup
    E_{\mathrm{build}}.
\]
这些错误不应混成一团。CLI 参数错误应在运行前报告；I/O 格式错误应指出记录边界或维度问题；
CUDA runtime 错误应保留底层诊断；pipeline 错误应说明并发阶段、输出缺包或线程异常；
构建错误应停留在 CMake 或编译器层。

kernel 层尤其需要错误翻译（error translation）。CUDA API 返回 \texttt{cudaError\_t}，
这是 runtime 语言；应用层更需要知道的是“哪次设备操作失败、失败码是什么、发生在何处”。
\texttt{JACOBI\_CUDA\_CHECK} 与 \texttt{CudaError} 的价值就在这里：把底层错误提升为稳定异常语义，
同时保留表达式、文件、行号和 CUDA 错误字符串。

还要区分两类常被混淆的失败：
\[
    \text{algorithmic non-convergence}
    \ne
    \text{CUDA runtime failure}.
\]
达到 \texttt{max\_sweeps} 仍未完全收敛，是数值算法状态；\texttt{cudaMalloc}、
kernel launch 或 \texttt{cudaMemcpy} 失败，是运行时系统状态。前者适合进入计算报告或验证指标，
后者通常应终止当前任务并传播诊断。把二者都叫“kernel failed”会遮蔽真正原因。

pipeline 层则负责跨线程错误传播。工作任务异常通过 \texttt{future.get()} 回到主线程；
写线程异常保存在输出阶段的异常指针中，并通过 \texttt{submit} 或关闭路径重抛。
因此，只要任一线程认为本次输出不可信，\texttt{JacobiSvdPipeline::run()} 就不能报告成功。
这就是 fail-fast 语义：宁可明确失败，也不要产出静默污染的实验文件。

\subsection{数据结构边界：同名“矩阵”不是同一种东西}

从数据结构角度看，系统里至少存在六种不同层次的“矩阵”：

\begin{itemize}
    \item 文件中的矩阵：网络字节序头部、载荷和持久化布局；
    \item 派发任务中的矩阵：带序号的原始 payload 与 pinned host buffer；
    \item 主机侧矩阵：\texttt{io::Matrix}，即维度加行主序 \texttt{std::vector}；
    \item 设备侧矩阵：\texttt{DeviceMatrix}，即维度加 \texttt{double*}；
    \item 内部工作矩阵：可选的列主序 \texttt{d\_a\_layout}；
    \item 输出矩阵：按 \(U\)、\(1\times n\) 的 \(\Sigma\)、\(V\) 三张矩阵写入流。
\end{itemize}

它们不能混为一谈。文件矩阵适合 I/O 层，派发任务适合 pipeline 的输入边界，
\texttt{io::Matrix} 适合应用层和测试；\texttt{DeviceMatrix} 适合 kernel wrapper。
列主序工作区只属于 \texttt{kernels.cu} 内部。若把设备指针暴露给 pipeline，
应用层就会开始关心 \texttt{cudaFree}；若把文件映射地址暴露给 kernel，
数值核心就会开始关心平台页缓存。两种方向都是坏味道。

正确的数据流是逐层降噪：
\[
    \text{file bytes}
    \xrightarrow{\text{decode or dispatch}}
    \text{host matrix}
    \xrightarrow{\text{copy}}
    \text{device matrix}
    \xrightarrow{\text{compute}}
    \text{host result}
    \xrightarrow{\text{encode}}
    \text{file bytes}.
\]
其中 \texttt{*.mat} 高吞吐路径把第一段细化为：
\[
    \text{file bytes}
    \xrightarrow{\text{single cursor}}
    \texttt{MatDispatchTask}
    \xrightarrow{\text{worker decode}}
    \texttt{io::Matrix}.
\]
复杂性没有消失，但它被放到了可定位的位置。

\subsection{良构系统的不变量}

围绕纯函数式 kernel，可以列出本项目希望长期保持的一组系统不变量。

\paragraph{不变量一：领域核心只依赖内存中的类型化数据。}
kernel 入口不接收路径、文件句柄、原始字节流或 CLI 字符串。所有外部输入必须先被解析为维度明确、
元素类型明确、布局明确的矩阵对象。

\paragraph{不变量二：公开结果不暴露设备资源。}
\texttt{JacobiSvdResult} 返回主机侧值，不返回 device pointer、stream、event 或临时缓冲区。
上层可以移动、序列化和测试结果，但不能管理 kernel 内部资源。

\paragraph{不变量三：应用层只编排，不重写领域逻辑。}
pipeline 可以决定何时读、何时调优、何时提交线程池、何时写出和如何处理背压；
但不能复制 Jacobi 旋转公式，也不能手写 SVD 后处理逻辑。

\paragraph{不变量四：基础设施层不决定算法语义。}
I/O 层可以选择 \texttt{*.mat} 或文本格式，可以使用 mmap、普通流或 pinned host buffer；
但不能改变 \(\epsilon\)、最大 sweep 数、列对调度或奇异值语义。

\paragraph{不变量五：表示层不越过应用层。}
CLI 只能生成配置并调用 pipeline。任何为了方便而在 \texttt{main.cpp} 中直接读文件、
写矩阵或调用内部 kernel helper 的做法，都应视为架构回归。

\paragraph{不变量六：错误在边界被翻译，而不是在深处格式化给用户。}
底层错误应携带诊断信息向上传播，由合适层决定用户可见表达。
kernel 不打印用户帮助，I/O 不决定进程退出码，pipeline 不解析命令行。

\paragraph{不变量七：并发不改变输出契约。}
线程池可以让 testcase 乱序完成，但输出文件仍必须按输入顺序写出
\((U_i,\Sigma_i,V_i)\)。重排序缓冲区是应用层机制，不是数学层机制。

\paragraph{不变量八：构建 artifact 不改变运行语义。}
CMake 可以选择 nvcc driver mode 或 CMake CUDA language 路径；同一输入、同一配置和同一源码
应暴露相同应用层契约。性能可能变化，语义不应变化。

\subsection{从模块图到依赖方向}

把上述思想压缩为依赖图，可以得到：
\[
    \begin{array}{ccccc}
        \texttt{main.cpp}
         & \longrightarrow &
        \texttt{pipeline.cpp}
         & \longrightarrow &
        \{\texttt{io.cu},\texttt{kernels.cu}\}.
    \end{array}
\]
这张图有两个关键含义。

第一，\texttt{io.cu} 与 \texttt{kernels.cu} 在应用层之下并列。
它们分别处理数据进入系统和数值核心计算，二者不应相互依赖。
kernel 不读取文件；I/O 不执行 SVD。pipeline 是把它们组合成用例的地方。

第二，依赖方向从用户意图到应用用例，再到具体能力。外层知道内层，内层不反向知道外层。
越靠近数学核心，变化频率应越低；越靠近用户和平台，变化频率可以越高。
可以用更抽象的公式表达：
\[
    \text{System}
    =
    \text{Adapters}
    \circ
    \text{Application Use Case}
    \circ
    \text{Domain Kernel}.
\]
这里的 \(\circ\) 不是单纯函数复合，而是责任边界的复合：
每层只把下一层需要的最小信息传过去。

\subsection{当前实现的边界风险}

当前实现已经体现了上述方向，但仍有几处需要持续警惕。

\paragraph{第一，\texttt{JacobiSvdConfig} 同时包含数值参数与执行策略。}
这在当前规模下可以接受，但未来若加入 device id、stream policy、workspace 预算、profiling 开关，
最好区分数学配置、执行配置和实验观测配置。否则配置对象会变成“所有旋钮的大口袋”。

\paragraph{第二，kernel 层仍是同步接口。}
\texttt{one\_sided\_jacobi\_svd} 当前是同步风格：调用完成后结果已回到主机。
pipeline 可以跨 testcase 并发提交 host 任务，但这不等价于显式 CUDA stream 级流水线。
若未来要实现 host decode、host-device copy、kernel、device-host copy 的真正重叠，
需要新增异步 compute port，而不是在 pipeline 里偷偷管理 stream 生命周期。

\paragraph{第三，执行配置仍有正确性含义。}
\texttt{threads\_per\_block} 看似只是性能参数，但上一章已经指出，当前共享内存二分规约对线程数形态有隐含假设。
默认 \(256\) 是安全点；后续最好把合法值收紧为 \(2\) 的幂，或修改规约以正确处理任意 warp 对齐线程数。
性能旋钮不能悄悄参与正确性。

\paragraph{第四，结果报告仍较粗。}
kernel 只报告 \texttt{sweeps}；pipeline 报告 testcase 数、写出矩阵数、总 sweep 与布局调优摘要。
这足以表达语义成功，但还不是性能遥测（performance telemetry）。
若要定位吞吐瓶颈，应新增独立 profiling 结构，记录阶段耗时、队列等待、重排序峰值与 CUDA 计时。
不要把稳定报告对象变成调试日志。

\paragraph{第五，pinned buffer 是资源优化，不是无限内存许可证。}
\texttt{PinnedHostTaskBuffer} 让 \texttt{*.mat} 任务更适合派发，但页锁定内存会影响操作系统调页自由度。
真正的上界来自 pipeline 的 in-flight 窗口和完成队列背压，而不是 pinned allocation 本身。
如果窗口无限增大，高吞吐路径也会变成资源耗尽路径。

\paragraph{第六，边界不应被性能优化轻易打穿。}
高性能项目常有一种诱惑：为了少拷贝、少封装、少函数调用，把文件映射地址、pinned buffer、
device pointer、CUDA stream 和输出队列揉成一个大对象。除非有明确测量证据，并且仍能保持语义边界，
否则这种优化更像债务而不是性能。好优化应该让特殊情况变成局部策略，而不是让特殊情况污染整个系统。

\subsection{小结}

本章把前一章的 CUDA kernel 从“实现细节”提升为“系统构件”来观察。它应当像纯函数一样被调用：
\[
    (A,C)\mapsto(U,\Sigma,V,s),
\]
即使其内部不可避免地包含设备内存、kernel launch、布局转置、runtime 同步和错误检查。
这个函数式外观不是美学装饰，而是可测试性、可替换性、错误隔离和长期演化的基础。

轻量 DDD 在这里的价值，不是给项目贴企业架构标签，而是提醒我们：真正的领域核心是 SVD 的数值不变量
和 CUDA 执行约束；CLI、I/O、pipeline 与 CMake 都是围绕它建立的边界。
kernel 不认识文件，不认识路径，不认识输出顺序，不认识用户脚本，也不认识构建目录。
它只认识矩阵、配置和数值结果。

这也让后续章节自然衔接：下一章讨论 I/O 如何把外部字节转化为矩阵流与 pinned dispatch task；
再下一章讨论 pipeline 如何用全局线程池、future 窗口和重排序缓冲区组合输入流、纯函数式 kernel 与输出流；
随后 CLI 与 CMake 分别定义用户意图边界和构建 artifact 边界。
一个良构 CUDA 数值系统的脊梁就在这里：把最复杂、最昂贵、最脆弱的 kernel 放在最窄、最稳定、最可验证的位置。
